% Generated draft: requires mathematical and source review.
\documentclass[11pt]{article}
\usepackage[margin=25mm]{geometry}
\usepackage{fontspec}
\setmainfont{Latin Modern Roman}
\usepackage{amsmath,amssymb,mathtools,unicode-math}
\setmathfont{Latin Modern Math}
\usepackage{enumitem,needspace,xurl,hyperref}
\setlength{\parindent}{0pt}
\setlength{\parskip}{5pt}
\setlength{\emergencystretch}{3em}
\newcommand{\N}{\mathbb{N}}
\newcommand{\Z}{\mathbb{Z}}
\newcommand{\Q}{\mathbb{Q}}
\newcommand{\R}{\mathbb{R}}
\newcommand{\C}{\mathbb{C}}
\newcommand{\F}{\mathbb{F}}
\newcommand{\A}{\mathbb{A}}
\newcommand{\PP}{\mathbb{P}}
\newcommand{\E}{\mathbb{E}}
\newcommand{\eps}{\varepsilon}
\newcommand{\dd}{\,\mathrm{d}}
\begin{document}
\section*{UnsolvedMath ID: 30003200}
\textbf{Source problem code:} catalogue entry
\subsection{Definitions and mathematical statement}
A locally finite tree is a connected graph without cycles in which every vertex has finite degree. Its boundary \(\partial T\) is the set of ends: equivalence classes of infinite simple rays, where two rays are equivalent if they eventually share a common tail. Give \(\operatorname{Aut}(T)\) the topology of pointwise convergence on the discrete vertex set. A closed subgroup \(G\leq\operatorname{Aut}(T)\), with the induced topology, is a locally compact group. If \(\partial T\neq\varnothing\), its boundary action is transitive when
\[
\forall\xi,\eta\in\partial T\quad\exists g\in G\quad g\xi=\eta.
\]
A locally compact group \(G\) is of type I if every nonzero strongly continuous unitary representation \(\pi:G\to\mathcal U(H)\) whose generated von Neumann algebra \(\pi(G)''\) has center \(\mathbb C I_H\) generates a type I factor; that is, \(\pi(G)''\) is isomorphic, as a von Neumann algebra, to \(\mathcal B(K)\) for some complex Hilbert space \(K\).

The supplied question asks: for every locally finite tree \(T\) and every closed subgroup \(G\leq\operatorname{Aut}(T)\) acting transitively on \(\partial T\), must \(G\) be of type I?

The source states no infinitude assumption on \(T\), no lower bound on the number of ends, and no additional hypothesis on \(G\). It does not specify whether an action on an empty boundary counts as transitive. These qualifications remain unresolved; the supplied formulation is retained without adding restrictions. Its identification with the intended folklore conjecture requires human review.
\subsection{Short English statement}
The question asks whether every closed group of automorphisms of a locally finite tree that acts transitively on its ends is a type I group. The supplied record labels the problem partially solved, but the intended treatment of trees with few or no ends and any missing hypotheses require review.
\subsection{Sources}
\begin{itemize}
\item Local UnsolvedMath registry, supplied record with permanent ID \(30003200\), source problem code \texttt{catalogue entry}, title ``Type-I Rigidity of Boundary-Transitive Tree Groups,'' in the Oberwolfach Reports Open Problems collection. The record gives proposed year \(2016\) and names no individual proposer.
\item Supplied report citation: \(C^*\)-Algebras (2017), Oberwolfach Reports, pp. 2269--2345, DOI \texttt{10.4171/owr/2016/40}. \url{https://doi.org/10.4171/owr/2016/40}. The record describes these reports as documenting mathematical workshops and their open-problem discussions. The citation's stated year \(2017\) and the record's proposed year \(2016\) are both retained.
\item Supplied literature source: ``Houdayer and Raum Locally compact groups acting on trees.'' \url{https://doi.org/10.4171/CMH/458}.
\item Supplied literature source: ``Caprace Kalantar and Monod Type I conjecture.'' \url{https://arxiv.org/abs/2110.00190}.
\item Literature assessment dated exactly \texttt{2026-08-21}: the supplied curation labels the problem ``partially solved'' and states that the type I conjecture is proved for major classes, including the Burger Mozes setting, while the full boundary transitive tree group classification remains open. This is supplied dated curation derived from the OWR record, not an independently verified assessment. Its time-sensitive status and cited claims require independent verification.
\item The supplied statement status is \texttt{corrected\_verified}. Its verification note says that the official report was checked and identifies the statement as its folklore type I conjecture with the boundary action made explicit. This is a supplied verification claim; no independent source visit or verification is asserted here.
\end{itemize}
\end{document}
